\chapter{Representations and Path Algebras}
\label{ch:rm03-quivers}

Let \(k\) be a field. Algebras and module actions are unital, and
\(\operatorname{End}_k(V)\) is multiplied by composition. We use left
modules throughout. In particular, an arrow from \(i\) to \(j\) will
act from the vertex summand at \(i\) to the vertex summand at \(j\).

\section{Representations of an algebra}

\begin{proposition}[RM03-C06.01. Algebra representations and modules]
\label{prop:rm03-algebra-representations}
For a \(k\)-algebra \(A\) and a \(k\)-vector space \(V\), unital
\(k\)-algebra homomorphisms
\(\rho\colon A\to\operatorname{End}_k(V)\) correspond bijectively to
left \(A\)-module structures on \(V\) extending its given
\(k\)-vector-space structure. A linear map intertwines two such
representations precisely when it is \(A\)-linear.

If \(I\) is a two-sided ideal of \(A\), restriction along
\(A\to A/I\) identifies left \(A/I\)-modules with left \(A\)-modules
annihilated by \(I\), with all module maps preserved.
\end{proposition}

\begin{proof}
Set \(av=\rho(a)(v)\). The linearity of \(\rho\) and of the operators
gives the distributive laws and the given scalar action.
Multiplicativity gives \((ab)v=a(bv)\), and \(\rho(1_A)=1_V\) gives
the unit law. In the other direction \(a\) acts by a linear map
because \(a(\lambda v)=(a(\lambda 1_A))v=\lambda(av)\), and the
action laws give the required \(k\)-algebra homomorphism. The two
constructions recover the same operator at every \(a\). The
intertwining condition says exactly that \(f(av)=af(v)\).

An action obtained by restriction is zero on \(I\). Conversely, if
\(IV=0\), the formula \((a+I)v=av\) is independent of the
representative and satisfies the module laws. A map commutes with the
action of \(A/I\) exactly when it commutes with the action of all
representatives in \(A\).
\end{proof}

The group case starts with \(G\to\operatorname{GL}_k(V)\), and the
linear extension to \(k[G]\) gives the action in
Proposition~\ref{prop:rm03-algebra-representations}. The group algebra
correspondence, Schur's lemma, and Maschke's theorem with its
characteristic hypothesis are established in
\cite[Ch.~8]{p103}. In the quiver case, the generators act between
specified summands, and the multiplication of paths controls which
composites are defined.

\section{Paths, idempotents, and vertex spaces}

\begin{definition}[RM03-C06.02. Quivers and path algebras]
\label{def:rm03-path-algebra}
A finite quiver \(Q=(Q_0,Q_1,s,t)\) has a finite nonempty set
\(Q_0\) of vertices, a finite set \(Q_1\) of arrows, and source and
target maps \(s,t\colon Q_1\to Q_0\). Parallel arrows and loops are
allowed. For each vertex \(i\) there is a path \(e_i\) of length zero
with source and target \(i\). A positive-length path is written
\(p=a_\ell\cdots a_1\), where \(a_1\) is traversed first and
\(t(a_b)=s(a_{b+1})\) for \(1\leq b<\ell\).

The path algebra \(kQ\) is the vector space with basis the paths.
For basis paths \(p,q\), the product \(pq\) is their concatenation
with \(q\) traversed first when \(s(p)=t(q)\), and is zero
otherwise. Concatenation with a length-zero path at a matching
vertex leaves the other path unchanged. Multiplication extends
bilinearly to finite linear combinations.
\end{definition}

\begin{lemma}[RM03-C06.03. Multiplication and the path basis]
\label{lem:rm03-path-basis}
The multiplication makes \(kQ\) an associative unital \(k\)-algebra.
The vertices give pairwise orthogonal idempotents with sum \(1\).
The space \(e_j(kQ)e_i\) has basis the paths from \(i\) to \(j\).
The algebra \(kQ\) is finite-dimensional exactly when \(Q\) has no
oriented cycle.
\end{lemma}

\begin{proof}
For paths \(p,q,r\), both \((pq)r\) and \(p(qr)\) vanish unless
\(s(p)=t(q)\) and \(s(q)=t(r)\). When both equalities hold, each
parenthesization yields the same ordered list of arrows, starting
with \(r\). Thus multiplication is associative on the basis and
hence on \(kQ\).

The length-zero paths satisfy \(e_j e_i=0\) for \(i\ne j\) and
\(e_i^2=e_i\). For a basis path \(p\), the product
\(e_j p e_i\) is \(p\) if \(s(p)=i\) and \(t(p)=j\), and is zero
otherwise. This proves the corner basis claim. Summing the left or
the right vertex action gives the identity on every path, so
\(\sum_i e_i=1\).

The powers of an oriented cycle have distinct lengths and are
linearly independent basis elements. If there is no oriented cycle,
no path visits a vertex twice. A repeated vertex would give a
positive-length closed path, and a shortest closed subpath would
give an oriented cycle. All lengths are then at most
\(\lvert Q_0\rvert-1\). Finiteness of the arrow set gives finitely
many paths of these lengths.
\end{proof}

\begin{definition}[RM03-C06.04. Quiver representations]
\label{def:rm03-quiver-representations}
A representation \(V\) of \(Q\) consists of vector spaces \(V_i\)
at the vertices and linear maps \(V_a\colon V_{s(a)}\to V_{t(a)}\)
on the arrows. For a path \(p=a_\ell\cdots a_1\) put
\(V_p=V_{a_\ell}\circ\cdots\circ V_{a_1}\), and set
\(V_{e_i}=1_{V_i}\). A morphism \(f\colon V\to W\) is a family of
linear maps \(f_i\colon V_i\to W_i\) satisfying
\(f_{t(a)}V_a=W_a f_{s(a)}\) for each \(a\). Composition is taken
at each vertex.
\end{definition}

\begin{theorem}[RM03-C06.05. Quiver representations and path algebra modules]
\label{thm:rm03-quiver-equivalence}
The category of left \(kQ\)-modules is \(k\)-linearly equivalent to
the category of representations of \(Q\). The equivalence sends a
module \(M\) to the spaces \(e_iM\) with \(a\colon i\to j\) acting
by \(m\mapsto am\) from \(e_iM\) to \(e_jM\). Its inverse sends
\(V\) to \(\bigoplus_i V_i\). The equivalence restricts to
finite-dimensional objects and identifies module exactness with
exactness at every vertex.
\end{theorem}

\begin{proof}
Since the idempotents sum to \(1\), every \(m\in M\) is a sum of
the \(e_i m\). Applying \(e_j\) to any relation among these
summands isolates the summand at \(j\), proving that the sum is
direct. By Lemma~\ref{lem:rm03-path-basis}, \(e_j a e_i=a\) for
an arrow \(a\colon i\to j\), which gives the indicated arrow map.
A module map restricts to the summands and commutes with these
arrow maps.

Conversely, on \(\bigoplus_i V_i\) let \(p\colon i\to j\) act by
projecting to \(V_i\), applying \(V_p\), and including into
\(V_j\). The operators for composable paths multiply as the
operator of their concatenation. For noncomposable paths the
intermediate inclusion followed by projection is zero. The sum of
the vertex operators is the identity. Extending the operators
linearly therefore gives a unital \(k\)-algebra homomorphism, so
Proposition~\ref{prop:rm03-algebra-representations} gives the
module. The direct sum of the components of a morphism commutes
with the path operators, by induction on path length using the
arrow condition, and hence is a module map.

For a module, the canonical sum map \(\bigoplus_i e_iM\to M\)
commutes with the vertices and the arrows. Every basis path is a
product of these generators, so the sum map is a module
isomorphism. Starting with \(V\), each \(e_i\) becomes the
projection onto \(V_i\). Recovering those summands identifies
their arrow maps with the original \(V_a\). Both identifications
commute with morphisms by their formulas.

Each module map is block diagonal on the vertex summands, and its
kernel and image are direct sums of the component kernels and
images. Thus equality of image and kernel in a sequence is
equivalent to equality at each vertex. For arbitrary quiver
morphisms the vertex kernels are preserved by the arrow maps and
the arrows descend to the vertex quotients by the images, by the
same commutation law. This gives kernels and cokernels with the
universal properties at each vertex. Finally, the vertex set is
finite, so \(\bigoplus_i V_i\) is finite-dimensional if and only if
each \(V_i\) is finite-dimensional.
\end{proof}

\section{Relations and finite-dimensional quotients}

\begin{definition}[RM03-C06.06. Relations and admissible ideals]
\label{def:rm03-quiver-relations}
A relation from \(i\) to \(j\) is a finite linear combination
\(r=\sum_p c_p p\) of paths from \(i\) to \(j\). A representation
satisfies \(r=0\) if \(\sum_p c_p V_p\colon V_i\to V_j\) is zero.
Relations generate a two-sided ideal in \(kQ\). Every two-sided
ideal \(I\) can be described this way, since \(r\in I\) is the sum
of the components \(e_j r e_i\in I\).

Let \(J\) be the two-sided ideal generated by all the arrows.
An ideal \(I\) is admissible if there is an integer \(m\geq 2\)
such that \(J^m\subseteq I\subseteq J^2\).
\end{definition}

\begin{theorem}[RM03-C06.07. The module category with relations]
\label{thm:rm03-bound-quiver-equivalence}
Let \(I\) be generated as a two-sided ideal by a set of relations.
The equivalence in Theorem~\ref{thm:rm03-quiver-equivalence}
identifies left \(kQ/I\)-modules with quiver representations
satisfying those relations, with the same morphisms. Exactness in
this category is checked at each vertex.
\end{theorem}

\begin{proof}
Under the action on \(\bigoplus_i V_i\), the operator of a
relation \(r=\sum_p c_p p\) from \(i\) to \(j\) is the projection
to \(V_i\) followed by \(\sum_p c_p V_p\), followed by inclusion
of \(V_j\). Hence the relation acts by zero exactly when its
required component is zero. If the generating relations act by
zero, each term \(arb\), for \(a,b\in kQ\) and a generating
relation \(r\), also acts by zero. This is equivalent to
annihilation by \(I\). Proposition~\ref{prop:rm03-algebra-representations}
now gives the quotient action and all its module maps. Kernels,
images, and quotients of modules annihilated by \(I\) remain
annihilated by \(I\), so the exactness criterion follows from
Theorem~\ref{thm:rm03-quiver-equivalence}.
\end{proof}

\begin{proposition}[RM03-C06.08. Finiteness from admissibility]
\label{prop:rm03-admissible-finiteness}
For \(m\geq 1\), \(J^m\) is the span of all paths of length at
least \(m\). If \(J^m\subseteq I\), the quotient \(kQ/I\) is
finite-dimensional and is spanned by the images of the paths of
length below \(m\). If also \(I\subseteq J^2\), the images of all
length-zero and length-one paths are linearly independent.
\end{proposition}

\begin{proof}
The span of the positive-length paths is a two-sided ideal, and
every such path is a product of arrows, so it equals \(J\).
Every product of \(m\) positive-length paths vanishes or has
length at least \(m\). Conversely, a path of length at least \(m\)
can be cut into \(m\) composable positive-length paths. These two
observations identify the span with \(J^m\). There are finitely
many paths of length below \(m\), since \(Q_0\) and \(Q_1\) are
finite. The additional inclusion \(I\subseteq J^2\) says that a
linear combination of paths of length zero and one lying in \(I\)
lies in the span of paths of length two and above. Independence
of the path basis makes that combination zero.
\end{proof}

For the path convention, the module equivalence, and admissible
ideals, see \cite[Secs.~1.3--1.4]{quiver}. The two smallest
examples below also test which vertex idempotents occur on
either side of an arrow and how a relation changes its operator.

\begin{example}[RM03-C06.09. One arrow and a triangular matrix algebra]
\label{ex:rm03-one-arrow}
Take \(Q\) with vertices \(1,2\) and one arrow \(a\colon 1\to 2\).
The path basis is \((e_1,e_2,a)\), and the assignments
\[
 e_1\longmapsto E_{11},\qquad
 e_2\longmapsto E_{22},\qquad
 a\longmapsto E_{21}
\]
give an isomorphism to the \(2\times 2\) lower triangular
matrix algebra over \(k\), with \(E_{ij}\) the usual matrix
units. Indeed, \(E_{22}E_{21}E_{11}=E_{21}\),
\(E_{21}^2=0\), the other vertex multiplications on \(E_{21}\)
are zero, and the diagonal matrix units are orthogonal and
sum to the identity. These are the products of the path basis.
The three matrix units span the triangular algebra and are linearly
independent. Here \(J^2=0\), so \(I=0\) is admissible.
Explicitly, a single map \(T\colon V_1\to V_2\) gives the action
\[
 \begin{pmatrix}x&0\\ u&y\end{pmatrix}(v_1,v_2)
   =\bigl(xv_1,\;uT(v_1)+yv_2\bigr).
\]
Indeed the displayed actions of \(E_{11},E_{22},E_{21}\) are the
two projections and the operator of \(a\) from \(V_1\) to \(V_2\).
\end{example}

\begin{example}[RM03-C06.10. One loop with a square-zero relation]
\label{ex:rm03-square-zero-loop}
For one vertex and one loop \(a\), the paths are \(a^n\)
for \(n\geq 0\), with \(a^0=e\) the vertex. The map
\(x^n\mapsto a^n\) is multiplicative on a basis and hence gives
\(k[x]\simeq kQ\). The relation \(a^2=0\) generates \(I=J^2\),
so it is admissible and \(kQ/I\simeq k[x]/(x^2)\) has basis
the classes of \(1,x\). A module is a vector space \(V\) with a
linear operator \(T=V_a\) satisfying \(T^2=0\). A morphism to
another such module with operator \(U\) is precisely a linear
map \(f\) with \(fT=Uf\). These claims follow from
Theorem~\ref{thm:rm03-bound-quiver-equivalence}, and the explicit
action is \((\lambda+\mu x)v=\lambda v+\mu T(v)\).
\end{example}
